%%% Importing Packages

\usepackage{datetime} % Formatting dates.
\usepackage{fancyhdr} % Formatting headers.
\usepackage{lastpage} % Adding page numbers.

\usepackage{graphicx} % Inserting images.
\usepackage{float} % Postitioning images.

\usepackage{dirtytalk} % Speech and quote marks.
\usepackage{hyperref} % Hyperlinks and references.
\usepackage{xcolor} % Colouring text.
\usepackage{enumitem} % Formatting ordered lists. 

\usepackage{amssymb} % Additional maths symbols.
\usepackage{amsmath} % Additional maths tools.

\usepackage{algpseudocode} % Adding pseudocode.
\usepackage{algorithm} % Adding algorithms.
\usepackage{minted} % Adding code snippets.

\usepackage[margin=1in]{geometry} % Reduces the page margins.

%%% 



%%% Custom Formatting

\newcommand{\templateversion}{4.0.0}

\hypersetup{colorlinks=true, urlcolor=black, linkcolor=black} % Formatting URLs and links.

\newdateformat{monthyeardate} % The date format for the title page.
{%
    \monthname[\THEMONTH] \THEYEAR
}
\newdateformat{yeardate} % The date format for the copyright statement.
{%
    \THEYEAR
}

%%%



%%% Custom Commands

\newcommand{\colouredText}[2]{\color{#1}#2\color{black}}
\newcommand{\separate}[1]{\vspace{#1 pt}\noindent}

\newcommand{\comment}[1]{\colouredText{red}{#1}} % Allows you to leave yourself comments.

% Set notation.
\newcommand{\R}{\mathbb{R}}
\newcommand{\N}{\mathbb{N}}
\newcommand{\C}{\mathbb{C}}
\newcommand{\Q}{\mathbb{Q}}
\newcommand{\Z}{\mathbb{Z}}
\newcommand{\I}{\mathbb{I}}

% Maths notation.
\newcommand{\hessian}[1]{\mathbf{H}_{#1}}
\newcommand{\jacobian}[1]{\mathbf{J}_{#1}}
\newcommand{\lagrange}{\mathcal{L}}
\newcommand{\dotproduct}[2]{\langle\mathbf{#1},\mathbf{#2}\rangle}

% Join symbols.
\newcommand{\ojoin}{\setbox0=\hbox{$\bowtie$}%
    \rule[-.02ex]{.25em}{.4pt}\llap{\rule[\ht0]{.25em}{.4pt}}}
\newcommand{\insertinnerjoin}{\mathbin{\bowtie\mkern-5.8mu}\;}
\newcommand{\insertleftouterjoin}{\mathbin{\ojoin\mkern-5.8mu\bowtie}}
\newcommand{\insertrightouterjoin}{\mathbin{\bowtie\mkern-5.8mu\ojoin}}
\newcommand{\insertfullouterjoin}{\mathbin{\ojoin\mkern-5.8mu\bowtie\mkern-5.8mu\ojoin}}

\newcommand{\insertimage}[3]
{
    \begin{figure}[H]
        \centering
        \includegraphics[width=#3\linewidth]{#1}
        \caption{#2}
        \label{fig:#2}
    \end{figure}
}

\newcommand{\insertalgorithm}[4]
{
    \begin{algorithm}[H]
        \caption{#1}
        \begin{algorithmic}
            \State \textbf{Input}: #2
            \State \textbf{Output}: #3
            \separate{7}#4
        \end{algorithmic}
        \label{alg:#1}
    \end{algorithm}
}

\newcommand{\inserttable}[4]
{
    \begin{table}[H]
        \centering
        \begin{tabular}{|#1|}
            \hline #3 \\ \hline #4 \\ \hline
        \end{tabular}
        \caption{#2}
        \label{tab:#2}
    \end{table}
}

%%%



%%% Academic Environments

% Counters for automatic numbering.
\newcounter{definitionnumber}[subsubsection]
\newcounter{examplenumber}[subsubsection]
\newcounter{conjecturenumber}[subsubsection]
\newcounter{entrynumber}[subsubsection]
\newcounter{theoremnumber}[subsubsection]
\newcounter{lemmanumber}[subsubsection]
\newcounter{propositionnumber}[subsubsection]

% Custom enviroments.
\newcommand{\definition}[2]{\textbf{Definition \thesubsubsection.\number\numexpr\value{definitionnumber}+1\relax}: #2\stepcounter{definitionnumber}\label{def:#1}}
\newcommand{\example}[2]{\textbf{Example \thesubsubsection.\number\numexpr\value{examplenumber}+1\relax}: #2\stepcounter{examplenumber}\label{ex:#1}}
\newcommand{\conjecture}[2]{\textbf{Conjecture \thesubsubsection.\number\numexpr\value{conjecturenumber}+1\relax}: #2\stepcounter{conjecturenumber}\label{cnj:#1}}
\newcommand{\numberedentry}[3]{\textbf{#2 \thesubsubsection.\number\numexpr\value{entrynumber}+1\relax}: #3\stepcounter{entrynumber}\label{entry:#1}}
\newcommand{\proof}[2]{\par\separate{3}\textit{Proof}: #2 \begin{flushright}\(\square\)\end{flushright}\label{prf:#1}}
\newcommand{\solution}[2]{\par\separate{3}\textit{Solution}: #2 \begin{flushright}\(\square\)\end{flushright}\label{sol:#1}}
\newcommand{\theorem}[3]{\textbf{Theorem \thesubsubsection.\number\numexpr\value{theoremnumber}+1\relax}: #2 \proof{thm:#1}{#3}\stepcounter{theoremnumber}\label{thm:#1}}
\newcommand{\lemma}[3]{\textbf{Lemma \thesubsubsection.\number\numexpr\value{lemmanumber}+1\relax}: #2 \proof{lem:#1}{#3}\stepcounter{lemmanumber}\label{lem:#1}}
\newcommand{\proposition}[3]{\textbf{Proposition \thesubsubsection.\number\numexpr\value{propositionnumber}+1\relax}: #2 \proof{prop:#1}{#3}\stepcounter{propositionnumber}\label{prop:#1}}
\newcommand{\standardproblem}[3]{\textbf{Standard Problem \(\#\)#1}: #2 \solution{#1}{#3}\label{stpr:#1}}

\newcommand{\elemref}[2]{\hyperref[#1]{#2 \ref*{#1}}}
\newcommand{\preenvref}[3]{\hyperref[#1]{#2 \ref*{#1}.\number\numexpr\value{#3}+1\relax}}
\newcommand{\postenvref}[3]{\hyperref[#1]{#2 \ref*{#1}.\number\numexpr\value{#3}\relax}}

%%%



%%% Document Setup

% Title Page.
\title{\thetitle\\[0.5em]\large{\thesubtitle}}
\author{\theauthor}
\date{\thedate} 

% Custom Page Header.
\fancyhead{}
\fancyhead[R]{\thetitle\text{ - }\thesubtitle}

% Custom Page Footer.
\fancyfoot{}
\fancyfoot[L]{Written by \theauthor \\ 
\textit{\LaTeX\ File Template v.\templateversion: Copyright \date{\yeardate\today}, Matthew Emerson}} % Do not edit.
\fancyfoot[R]{Page \thepage \, of\, \pageref{mainDocumentEnd}\relax}

%%%